\documentclass[11pt]{article}
\usepackage[margin=1in]{geometry}
\usepackage[T1]{fontenc}
\usepackage[utf8]{inputenc}
\usepackage{lmodern}
\usepackage{microtype}
\usepackage[hidelinks]{hyperref}
\usepackage{fancyhdr}
\pagestyle{fancy}
\setlength{\headheight}{14pt}
\fancyhf{}
\fancyhead[L]{Working Paper}
\fancyhead[R]{The Capability Is in the Contest}
\fancyfoot[C]{\thepage}
\setlength{\parindent}{0pt}
\setlength{\parskip}{0.65em}
\title{\textbf{The Capability Is in the Contest}\\\large Prompt Battles and the Problem of Measuring Interactive AI}
\author{Watson Hartsoe\\\small Working Paper -- AI-augmented research process}
\date{17 August 2026}
\begin{document}
\maketitle
\begin{abstract}
Large-language-model evaluation usually speaks as if capability belonged to a portable object called the model. Interactive use makes that assumption increasingly difficult to sustain. A model that fails under one prompt may succeed under another; a skilled operator may uncover behavior unavailable to a novice; a public scoring rule can redirect contestant strategy; and a benchmark that changes its challenges in response to observed failures can become part of the causal system it claims to measure. This paper develops a stronger interpretation of Prompt Battles: not as a game for proving whether AI can or cannot do something, but as an adversarial elicitation ecology for studying how capability appears, disappears, transfers, and mutates across model, operator, prompt history, evaluation rule, and adaptive challenge distribution. The argument has a limit. If every relevant component is folded into a coupled system, capability becomes impossible to localize and benchmarks lose their comparative object. I therefore distinguish model-local capability, elicited capability, coupled capability, and ecological performance, and propose a two-track evaluation architecture: sealed anchors for comparability and frontier battles for discovery. The result is not a replacement for conventional benchmarks. It is a method for investigating the conditions under which conventional capability claims stop being well-formed.
\end{abstract}
\noindent\textbf{Keywords:} large language models; evaluation; capability elicitation; human-AI interaction; dynamic benchmarking; prompt sensitivity; adversarial collaboration; performative evaluation

\section{Introduction}
A Prompt Battle seems, at first, to pose an ordinary question in a theatrical form. One participant tries to make a model exhibit a familiar limitation; another tries to prompt past it. The apparent research object is the model. The battle is merely an energetic way to test it.

That interpretation collapses as soon as the battle works.

Suppose a defender elicits a biased answer and a challenger then obtains an unbiased answer by adding an explicit constraint. Nothing has been settled about whether the model "is biased." The two participants have established two compatible facts: there exists a route to failure, and there exists a route to success. Prompt-sensitivity work makes this more than a philosophical nuisance. Small, intent-preserving changes can alter model responses, so one completion is not a stable sample of a capability unless the evaluation specifies how prompts are distributed and varied \cite{chatterjee2024posix}. Behavioral-testing work such as CheckList makes the parallel point that aggregate accuracy can conceal systematic failures that only appear under targeted tests \cite{ribeiro2020beyond}.

The battle therefore exposes a defect in the proposition it was designed to adjudicate. "Can the model do X?" is under-specified. Does can mean that at least one prompt succeeds? That a typical user can elicit success? That an expert can? That success survives paraphrase? That it survives adversarial perturbation? That it remains available when the model's system instructions, tools, interface, and interaction history are held fixed? Each quantifier produces a different capability claim.

This paper argues that Prompt Battles become scientifically interesting at precisely this point of failure. Their strongest object is not a winner and not even a completion. It is an elicitation ecology: a coupled and partially adaptive system in which model behavior is generated by the relation among a model, an operator policy, a prompt history, an evaluation rule, and a challenge distribution. This reframing explains why several apparently separate methodological problems - prompt sensitivity, operator expertise, dynamic benchmarks, public scoring, judge disagreement, and adversarial collaboration - converge on one issue. The evaluation procedure is not outside the phenomenon. It helps constitute the behavior it later calls evidence.

That claim needs restraint. It would be easy to respond by declaring every tool, operator, policy, interface, and audience part of one enormous "system." Such a move explains everything and measures nothing. The paper therefore develops four levels of attribution and a corresponding research design. Model-local capability asks what survives across elicitation procedures. Elicited capability asks what can be recovered under a specified intervention budget. Coupled capability asks what a reliably available human-model configuration can do. Ecological performance asks what emerges in the full adaptive setting. The point is not to choose one level as the real one. The point is to stop sliding among them without noticing.

\section{From Binary Capability to Elicitation Depth}
The earliest Prompt Battle distinction between cannot, will not, and should not looks like a taxonomy of model behavior. At the interface, however, these categories can produce the same observation: the requested behavior does not occur. Failure is therefore not self-interpreting. A missing capability, a safety policy, a system instruction, a tool restriction, a weak prompt, and an evaluator who has not discovered the right intervention can all generate observationally similar traces.

The Elicitation Game makes this identification problem explicit by evaluating techniques that recover hidden capabilities from models whose relevant behavior is intentionally difficult to elicit \cite{hofstatter2025elicitation}. The lesson is not that every failure conceals a capability. It is that non-performance under a particular elicitation procedure does not identify incapacity. This changes the natural experimental unit. Instead of classifying a behavior as can or cannot, one can measure elicitation depth: the strongest intervention required to recover the behavior. Ordinary prompting, expert prompt search, multi-turn interaction, few-shot demonstration, system-instruction variation, tools, activation steering, or fine-tuning occupy different positions on such a ladder.

This also clarifies the logical error of a simple capture-the-flag design. Let B(M,p)=1 denote an operationalized failure of model M under prompt p. A defender who finds p with B(M,p)=1 and a challenger who finds q with B(M,q)=0 have not contradicted each other. They have sampled two regions of prompt space. A meaningful battle therefore needs to specify whether it evaluates existence, prevalence, robustness, worst-case behavior, typical-user accessibility, or expert recoverability. The winning condition cannot be chosen after the search.

This reframing makes prompt sensitivity part of the phenomenon rather than nuisance variance. The quantity of interest may be a distribution over behavior under a defined family of intent-preserving prompts, or a map of the routes by which behavior can be reached. A brilliant counter-prompt can then be interpreted in two opposed ways. As an existence proof, it demonstrates recoverability. As a robustness diagnostic, the very need for such a prompt may demonstrate fragility. The same completion changes meaning with the quantifier.

\section{The Operator Enters the Object}
Prompt Battles become even stranger when they distinguish the System Builder from the Prompt Pilot. The former establishes the governing sandbox and constraints; the latter chooses prompts within it. Once these roles are explicit, the completion can no longer be attributed to the model without qualification. The human chooses the next intervention based on previous outputs. In control terms, the prompt sequence is a policy over interaction history. Different operators implement different policies over the same model.

This is not merely a confound to eliminate. Kirsh and Maglio's distinction between pragmatic and epistemic action suggests why an expert operator may outperform a novice even when neither writes a single perfect instruction \cite{kirsh1994distinguishing}. Some actions are valuable because they alter the external problem state and make information easier to perceive or compute. The same possibility appears in prompting. A completion can be useless as a final answer yet valuable because it reveals a hidden distinction, produces a contrast, externalizes alternatives, or changes what the operator knows to ask next. The value of such a prompt is delayed and trajectory-dependent.

If that is correct, the prompt-completion pair is too small a unit for expert prompting. The operative object is a sequence in which intermediate outputs become temporary external states for thought. An expert advantage might therefore reside not in better first prompts but in a higher frequency of strategically placed epistemic prompts whose value becomes visible only in later turns.

This creates an attribution problem. If expert prompting is required for a behavior to appear, whose capability is it? Clark and Chalmers' extended-mind argument provides a useful pressure test without settling the issue \cite{clark1998extended}. Their account makes coupling, causal contribution, and reliable availability relevant to where a cognitive process is located. Applied cautiously to AI evaluation, it suggests a distinction between episodic assistance and stable coupled capability. A randomly encountered expert who rescues one task is not equivalent to an operator-model configuration that is reliably available, repeatedly integrated, and robust to interruption.

The distinction matters because practical AI systems are increasingly assembled from models, retrieval systems, tools, persistent histories, interfaces, and human operators. A benchmark that insists capability must belong to the bare model can misdescribe useful competence. A benchmark that attributes every successful configuration to "the system" can no longer say what changed when one component is removed. The empirical solution is not conceptual inflation but ablation and substitution: rotate operators across models, swap system builders, remove tools, interrupt histories, and measure which competences survive. Capability attribution should be an empirical dependency map.

\section{When the Benchmark Learns From the Model}
A second fracture appears when the critique pool adapts. Dynamic benchmarking already makes evaluation less static. Dynabench, for example, places humans and models in a loop where people construct challenge examples that expose current failures, feeding model development and subsequent assessment \cite{kiela2021dynabench}. Prompt Battles can go further because they may adapt not only examples but the critique itself. "The model is biased" may become "the model cannot recognize its bias" and then "the model cannot recover from adversarially induced bias." The task ontology moves in response to performance.

This can be productive. A fixed benchmark eventually saturates. A frontier evaluation that follows failures can continue discovering boundaries. But the properties that make it good at discovery make it bad as a stable measuring stick. If challenge selection at time t+1 depends on failures at time t, the evaluation distribution is endogenous.

Perdomo and colleagues' concept of performative prediction sharpens the feedback problem \cite{perdomo2020performative}. In their setting, deploying a prediction changes the distribution of future outcomes. A public Prompt Battle can be performative in an analogous, though not identical, sense. Published failures change human prompting strategies. Public rubrics change what contestants optimize. Model developers may train against visible failures. Safety policies may be patched. The critique pool adapts. The next evaluation then contains causal descendants of the previous one.

A longitudinal battle can therefore mistake evaluator-induced change for generalized technological progress. A model may improve on the visible frontier because the frontier helped teach the ecosystem what to optimize. The strongest design is consequently dual. A sealed anchor league preserves fixed prompts, criteria, and unpublished cases for longitudinal comparison. A frontier league adapts aggressively to failures and records the genealogy of every new challenge. The first estimates change against a stable reference. The second measures the rate and character of boundary discovery. Collapsing them into one score sacrifices both purposes.

\section{The Score Is an Intervention}
The temptation to rescue an unruly battle through better scoring creates a third problem. Early versions of the Prompt Battle archive respond to ambiguous winning by proposing standardized rubrics, public criteria, judges, crowdsourcing, and automated scoring. Later versions shift toward "judging the dance": depth, adaptability, innovation, synergy, and the evolving quality of the interaction. This is not merely a better rubric. It changes the theory of evidence from final-output adequacy to process quality.

Even a perfectly specified rubric does not remain outside the system. Manheim and Garrabrant's taxonomy of Goodhart effects is useful here because the contestants see the metric and receive direct selection pressure to optimize it \cite{manheim2018goodhart}. If creativity, depth, flag proximity, or prompting sophistication becomes the visible score, expert participants can search for behaviors that maximize those proxies. The fairer the game becomes in one sense - the more clearly its scoring is published - the easier it becomes to optimize the measurement surface rather than the intended construct.

This yields a structural conflict. Hidden criteria reduce strategic gaming but undermine procedural transparency. Public criteria increase fairness while increasing optimization pressure on the proxy. A practical design therefore needs held-out transfer measures. Contestants may train and compete under a public rubric, but winning strategies should then be tested on novel flags and by criteria not available during optimization. If leaderboard rank collapses under transfer, the battle has selected metric-specialized technique rather than the intended general capability.

Judge disagreement introduces a related but distinct problem. Treating all disagreement as noise is too simple, but treating all disagreement as pluralistic insight is equally weak. Hautli-Janisz, Schad, and Reed distinguish error, fuzziness, and ambiguity as different sources of disagreement in argument annotation \cite{hautli2022disagreement}. Prompt Battle judging should similarly preserve rationales before reconciliation. A disagreement that disappears after factual correction differs from one that tracks a fuzzy threshold, which differs again from one that survives shared facts because judges inhabit competing interpretations of the flag. The object is not disagreement itself but the mechanism that generated it.

\section{Adversaries Need a Way to Lose}
Competition alone does not make evidence severe. A defender assigned to prove a limitation will search failure space; a challenger assigned to defeat it will search success space. If neither side states in advance what evidence would weaken its position, the battle can become indefinitely renewable. Every apparent loss can be answered with another example.

Adversarial collaboration offers a stricter model. Clark and colleagues argue for disputants jointly designing tests of competing hypotheses rather than independently producing demonstrations favorable to their own view \cite{clark2022enemies}. The scientific work therefore begins before the contest: opponents must share ownership of a test dangerous to both positions. In a Prompt Battle, that means jointly specifying the flag, the relevant quantifier, what observations support it, what observations weaken it, and what outcomes would remain ambiguous before role assignment.

Yet even this does not close the problem. Mellers, Hertwig, and Kahneman's adversarial collaboration on conjunction effects produced jointly accepted experiments while leaving room for divergent interpretations of the results \cite{mellers2001frequency}. Agreement on procedure and observation does not guarantee agreement on theoretical consequence. The elasticity can migrate from experimental design into interpretation.

A stronger protocol therefore separates an outcome crux from an interpretation crux. Before the battle, each side should record not only predictions but how different possible outcomes would update its position. Afterward, revisions should remain allowed, because unexpected evidence can reveal distinctions nobody anticipated. But revisions should carry a receipt: what new evidence forced the change, what new distinction was introduced, and why it was absent from the original update rule. The goal is not to outlaw reinterpretation. It is to distinguish legitimate conceptual learning from post-hoc immunization.

\section{Four Levels of Capability Attribution}
The preceding problems converge on a simple requirement: every capability claim should declare its level of attribution.

Model-local capability concerns behavior that remains comparatively stable across a specified range of elicitation procedures, operators, histories, and interfaces. It is the closest object to conventional benchmark talk about "the model."

Elicited capability concerns behavior recoverable under a declared intervention budget. The claim is not that the model reliably does X, but that X can be reached with interventions of depth d under specified conditions.

Coupled capability concerns a reliably available configuration of model, operator, tools, and interaction resources. Its defining test is dependency: removing or substituting a component changes competence, and the coupling is stable enough to be operationally meaningful rather than episodic assistance.

Ecological performance concerns the full adaptive arena: public scores, changing challenge pools, social stakes, audiences, institutional rules, provider policies, and learning across rounds. Here the research question is not what capability belongs to a component but how behavior emerges from a system whose measurement procedure is itself active.

These levels should not be forced into one leaderboard. They answer different questions and require different controls. Their value is diagnostic. Many arguments about whether an AI "can" do something are not substantive disagreements at all; they are unacknowledged shifts between levels.

\section{What the Battle Metaphor Hides}
The elicitation-ecology account also places limits on the battle metaphor itself. Consalvo's argument that there is no sealed magic circle warns against imagining that game rules create an autonomous experimental world \cite{consalvo2009magic}. Provider policies, interface constraints, model updates, audience expectations, participant histories, and institutional stakes continue to operate inside the supposed arena. Boundary leakage is therefore not automatically contamination. It may reveal which outside variables constitute the phenomenon.

The archive's borrowing of the breaking cypher requires even greater care. Johnson's work treats the cypher as a historically and culturally situated practice shaped by Africanist aesthetics, ritual, embodied and aural-kinesthetic relations, originality, social difference, spirit, and liberation \cite{johnson2023dark,johnson2011bboying}. Calling a shared chat a cypher because turns circulate through a common space preserves the easiest feature to extract while risking erasure of the relations that make the form meaningful. Even an "operational genealogy" can repeat the problem if it translates only what an interface can formalize.

This is more than a terminological caution. It generalizes to AI method design. Technical operationalization selects what can become a variable, state, rule, or score. Features that resist formalization are easy to demote as peripheral. But the inability of an apparatus to represent something is not evidence that the thing does not matter. A Prompt Battle that borrows social forms should therefore maintain a translation ledger recording what was preserved, transformed, absent, unrepresentable, or deliberately excluded. Formalizability is a property of the apparatus, not a measure of cultural importance.

\section{A Two-Track Research Architecture}
The resulting methodology is deliberately asymmetric.

The sealed anchor track is boring by design. It freezes task definitions, quantifiers, prompt families, system configurations, held-out transfer measures, and adjudication procedures. It supports repeated comparison and ablation. Its job is to say what changed under stable conditions.

The frontier battle track is allowed to be alive. Participants search prompt space adversarially. Intermediate prompts may function as epistemic actions. New critiques may descend from discovered failures. Judges preserve rationales and disagreement mechanisms. Public scoring can shape strategy, so transfer tests remain outside the optimization loop. Every new challenge records its parent observation. Its job is not to produce a timeless score but to find the next question.

The tracks communicate but do not collapse. Frontier discoveries can nominate future anchor tests, but a versioned anchor is never silently rewritten. Anchor failures can seed frontier inquiry, but the exploratory descendants are marked as a new construct rather than a harder instance of the old one. This preserves both memory and motion.

The design also changes the meaning of winning. A participant can win a round without settling a capability claim. The research success criterion is stronger: did the battle expose a discriminating variable, a new elicitation path, a dependency, a failure of transfer, a disagreement mechanism, or a revised question that can itself be tested? The event remains competitive, but the archive treats victory as evidence only when the conditions of inference are explicit.

\section{Limitations and Unresolved Territory}
This argument is assembled from research traditions that do not naturally belong to one another. Prompt sensitivity and capability elicitation directly concern language-model evaluation. Extended cognition concerns the location of cognitive processes. Performative prediction concerns prediction systems that alter target distributions. Goodhart's law concerns proxy optimization. Perspectivist annotation concerns disagreement. Adversarial collaboration concerns scientific disputes. Game and cypher scholarship concern social and cultural practices. Their conjunction is a wager, not an inherited unified theory.

Several claims therefore remain tests rather than conclusions. Whether expert prompting is substantially composed of epistemic actions requires trajectory-level evidence. Whether stable human-model coupling deserves capability attribution requires operational criteria that may differ by domain. Whether public Prompt Battles materially alter later model behavior requires causal lineage that may be impossible to recover from proprietary development. Whether judge disagreement contains stable interpretive structure requires rationale data collected before consensus procedures erase it. Whether battle-derived prompting skill transfers beyond a public rubric requires held-out evaluation.

The deepest unresolved issue is ontological. Evaluation needs an object. Interactive AI use keeps making that object larger. If the object expands until it includes every causal influence, comparison becomes vacuous. If it shrinks to the model weights alone, evaluation may cease to describe the system people actually use. The four-level scheme is therefore a temporary discipline, not a final ontology. Its purpose is to force capability claims to declare what they belong to.

\section{Conclusion}
Prompt Battles become less interesting when they are treated as theatrical benchmarks and more interesting when they are allowed to break the benchmark concept open. A failure under one prompt and success under another do not settle a binary capability claim. They reveal a search space. A skilled Prompt Pilot does not merely provide a better input; the operator may restructure the cognitive and computational problem over a trajectory. A dynamic challenge pool does not merely keep evaluation fresh; it can make the benchmark endogenous and eventually performative. A scoring rubric does not merely observe contestants; it changes what they optimize. A panel does not merely add noise when it disagrees; its disagreements can arise from distinct mechanisms. Adversaries do not produce severe evidence merely by opposing one another; they need shared conditions under which their own side can lose.

The resulting object is not "the AI" in isolation. Neither is it an undifferentiated human-machine assemblage. It is a set of nested attribution levels connected by experiments that remove, substitute, vary, and seal components. Conventional benchmarks remain necessary for comparison. Frontier battles are useful for a different reason: they discover when the assumptions that make conventional comparison possible have begun to fail.

That is the wager. Prompt Battles should not ask, simply, whether the model can capture the flag. They should ask what had to become part of the system before the flag became capturable - and whether that system is the same one we thought we were measuring.

\section*{Acknowledgment of Process}
This working paper was assembled from a preserved zettel field. The paper is a current wager, not evidence. Exact card payloads, sources, prompts, graph relations, making history, and the assembly instrument are preserved in the accompanying SLIPCASE package.

\appendix
\section{Assembly Instrument}
The exact assembly instrument is preserved at \texttt{\_PROMPTS/SLIPCASE\_v15.55-AM\_\_ASSEMBLY\_INSTRUMENT.txt}. The package manifest records its hash.

\section{Making History}
The package records PROVIDED, PRESERVED, RETRIEVED, DERIVED, VERIFIED, and UNVERIFIED states in \texttt{000\_\_MAKING\_HISTORY.txt}.

\section{Replication Path}
Open \texttt{index.html} or \texttt{000\_\_REBUILD.txt}. Root zettel TXT cards are immutable evidence payloads; derived views can be regenerated from those cards plus the relation compiler and resource registry.

\begin{thebibliography}{99}
\bibitem{chatterjee2024posix} Chatterjee, Anwoy and Renduchintala, H S V N S Kowndinya and Bhatia, Sumit and Chakraborty, Tanmoy. POSIX: A Prompt Sensitivity Index For Large Language Models. 2024.
\bibitem{ribeiro2020beyond} Ribeiro, Marco Tulio and Wu, Tongshuang and Guestrin, Carlos and Singh, Sameer. Beyond Accuracy: Behavioral Testing of NLP Models with CheckList. 2020. DOI: 10.18653/v1/2020.acl-main.442.
\bibitem{hofstatter2025elicitation} Hofst{\"a}tter, Felix and van der Weij, Teun and Teoh, Jayden and Djoneva, Rada and Bartsch, Henning and Ward, Francis Rhys. The Elicitation Game: Evaluating Capability Elicitation Techniques. 2025. DOI: 10.48550/arXiv.2502.02180.
\bibitem{kirsh1994distinguishing} Kirsh, David and Maglio, Paul. On Distinguishing Epistemic from Pragmatic Action. 1994. DOI: 10.1207/s15516709cog1804\_1.
\bibitem{clark1998extended} Clark, Andy and Chalmers, David J.. The Extended Mind. 1998. DOI: 10.1093/analys/58.1.7.
\bibitem{kiela2021dynabench} Kiela, Douwe and Bartolo, Max and Nie, Yixin and Kaushik, Divyansh and Geiger, Atticus and Wu, Zhengxuan and Vidgen, Bertie and Prasad, Grusha and Singh, Amanpreet and Ringshia, Pratik and Ma, Zhiyi and Thrush, Tristan and Riedel, Sebastian and Waseem, Zeerak and Stenetorp, Pontus and Jia, Robin and Bansal, Mohit and Potts, Christopher and Williams, Adina. Dynabench: Rethinking Benchmarking in NLP. 2021.
\bibitem{perdomo2020performative} Perdomo, Juan and Zrnic, Tijana and Mendler-D{\"u}nner, Celestine and Hardt, Moritz. Performative Prediction. 2020.
\bibitem{manheim2018goodhart} Manheim, David and Garrabrant, Scott. Categorizing Variants of Goodhart's Law. 2018. DOI: 10.48550/arXiv.1803.04585.
\bibitem{hautli2022disagreement} Hautli-Janisz, Annette and Schad, Ella and Reed, Chris. Disagreement Space in Argument Analysis. 2022.
\bibitem{clark2022enemies} Clark, Cory J. and Costello, Thomas H. and Mitchell, Gregory and Tetlock, Philip E.. Keep Your Enemies Close: Adversarial Collaborations Will Improve Behavioral Science. 2022. DOI: 10.1037/mac0000004.
\bibitem{mellers2001frequency} Mellers, Barbara and Hertwig, Ralph and Kahneman, Daniel. Do Frequency Representations Eliminate Conjunction Effects? An Exercise in Adversarial Collaboration. 2001. DOI: 10.1111/1467-9280.00350.
\bibitem{consalvo2009magic} Consalvo, Mia. There is No Magic Circle. 2009. DOI: 10.1177/1555412009343575.
\bibitem{johnson2023dark} Johnson, Imani Kai. Dark Matter in Breaking Cyphers: The Life of Africanist Aesthetics in Global Hip Hop. 2023.
\bibitem{johnson2011bboying} Johnson, Imani Kai. B-Boying and Battling in a Global Context: The Discursive Life of Difference in Hip Hop Dance. 2011.
\end{thebibliography}
\end{document}
